\documentclass[11pt]{article}
\usepackage{fontspec,amsmath,amssymb,amsthm,hyperref}
\usepackage[margin=1in]{geometry}
\setlength{\emergencystretch}{3em}
\title{Conditional scalar splitting recurrence and its strict endpoint}
\author{LangevinLean mathematical exposition}
\date{}
\begin{document}\maketitle
\noindent Source module: \texttt{ProximalLangevin/Recurrence.lean}. Compile this proof document with XeLaTeX or LuaLaTeX. The accompanying HTML places the complete Lean declarations beside every proof.


This file proves scalar consequences of an explicitly supplied recurrence. It does not supply the stochastic one-step law estimate for the paper’s coefficient. The independently proved PIPGLA law theorem uses the constructed coupling route.


\section{Correct coefficients and their signs}


Set

\[q=\frac{\lambda(1-\mu\gamma)}\gamma,\qquad c=\lambda\left(\frac{2\gamma LD}{N}+\lambda C\right).\]

Under the displayed nonnegative scalar premises, \(q,c\ge0\). When \(\gamma>0\) and \(\mu\gamma<1\), multiplying by positive denominators gives \(q<1\iff\lambda<\gamma/(1-\mu\gamma)\). The endpoint must be excluded from a geometric error floor. The subgradient forcing is \(\lambda^2C\), without an additional division by \(N\).



\noindent Lean declarations: \texttt{splittingRatio}, \texttt{splittingForcing}, \texttt{splittingRatio\_nonneg}, \texttt{splittingRatio\_lt\_one\_iff}, \texttt{splittingForcing\_nonneg}.

\section{Sum a supplied affine recurrence}


Assume \(e_{n+1}\le q e_n+c\). For \(q\ge0\), induction gives \(e_n\le q^ne_0+c\sum_{j<n}q^j\). If \(q<1\), the finite geometric identity gives \(e_n\le q^ne_0+c(1-q^n)/(1-q)\). If also \(c\ge0\), drop the nonpositive subtraction \(-cq^n/(1-q)\) to get \(e_n\le q^ne_0+c/(1-q)\). All three conclusions retain their recurrence hypothesis.



\noindent Lean declarations: \texttt{splitting\_recurrence\_sum}, \texttt{splitting\_recurrence\_exact}, \texttt{splitting\_recurrence\_bound}.

\section{Treat q = 1 and equal steps honestly}


At \(q=1\), induction gives the finite-horizon bound \(e_n\le e_0+nc\); there is no finite geometric floor. When \(\lambda=\gamma\) and the named denominators are nonzero, direct field algebra gives

\[\frac c{1-q}=\frac\gamma\mu\left(\frac{2LD}{N}+C\right).\]

This identity preserves the unscaled \(C\)-term and does not assert that an actual stochastic error satisfies the input recurrence.



\noindent Lean declarations: \texttt{splitting\_recurrence\_endpoint}, \texttt{splitting\_bias\_equal\_steps}.

\end{document}
